\chapter{Pneumothorax}

\section*{Definition}
Pneumothorax is the presence of air in the pleural space leading to partial or complete lung collapse. It is classified as primary spontaneous (no underlying lung disease), secondary spontaneous (underlying lung pathology), traumatic, or tension (with hemodynamic compromise).

\section*{What Happens in Your Body}
Air enters the pleural space via rupture of a subpleural bleb (spontaneous), through a chest wall defect (traumatic), or from alveolar-pleural fistula (secondary). Accumulation of intrapleural air equalizes or reverses the negative pleural pressure gradient, causing lung collapse and, in tension physiology, mediastinal shift with impaired venous return.

\section*{Causes}
\begin{itemize}
  \item \textbf{Primary spontaneous:} Rupture of apical subpleural blebs in tall, thin young adults (typically ages 20--40), smoking (90\% risk attributable), genetic factors (Birt-Hogg-Dub\'e syndrome, Marfan syndrome)
  \item \textbf{Secondary spontaneous:} COPD/emphysema (most common), cystic fibrosis, asthma, TB, \textit{Pneumocystis} pneumonia, lung cancer, endometriosis (catamenial pneumothorax)
  \item \textbf{Traumatic:} Penetrating trauma (stab, gunshot), blunt trauma (rib fracture), iatrogenic (central line, thoracentesis, mechanical ventilation, lung biopsy)
  \item \textbf{Tension:} Any pneumothorax with one-way valve mechanism allowing air entry but not exit
\end{itemize}

\section*{When to See a Doctor}
See your doctor if:
\begin{itemize}
  \item Sudden-onset ipsilateral pleuritic chest pain with acute dyspnea
  \item Hypotension, distended neck veins, tracheal deviation away from affected side (tension physiology)
  \item Oxygen saturation less than 90\% or inability to speak full sentences
  \item History of recurrent pneumothorax or known underlying lung disease
  \item High-risk activities (scuba diving, flying within 24 hours of onset)
\end{itemize}

\section*{Related Symptoms}
\begin{itemize}
  \item Sudden sharp pleuritic chest pain, dyspnea, decreased breath sounds, hyperresonance to percussion, tachypnea, tachycardia
\end{itemize}
\textbf{Important:} A small primary spontaneous pneumothorax (less than 15--20\% of hemithorax) may be minimally symptomatic, while a tension pneumothorax requires immediate needle decompression without waiting for imaging confirmation.

\section*{Self-Care / Home Management}
\begin{itemize}
  \item For small (less than 2 cm), asymptomatic primary spontaneous pneumothorax: observation with outpatient follow-up in 2--4 weeks
  \item Avoid air travel until pneumothorax is fully resolved on chest X-ray
  \item Avoid scuba diving permanently after any pneumothorax (high recurrence risk)
  \item Smoking cessation counseling (reduces recurrence risk by 2--3 fold)
  \item Restrict heavy lifting, Valsalva, and strenuous exercise during healing
\end{itemize}

\section*{How Your Doctor Will Evaluate You}
\begin{itemize}
  \item Chest X-ray (upright expiratory PA view most sensitive): visible visceral pleural line, absence of lung markings peripheral to the line
  \item CT chest (high-resolution) to identify blebs, bullae, or underlying lung pathology, especially in secondary pneumothorax
  \item Ultrasound (bedside): absence of lung sliding, presence of lung point sign
  \item Arterial blood gas to assess oxygenation and ventilation in symptomatic people
  \item ECG to differentiate from acute coronary syndrome (may show right axis deviation, decreased QRS amplitude)
\end{itemize}

\section*{Treatment Options}
\begin{itemize}
  \item Observation (small less than 2 cm, asymptomatic) or supplemental high-flow oxygen (accelerates reabsorption 4-fold)
  \item Needle aspiration (14--16 gauge, second intercostal space, midclavicular line) for symptomatic primary pneumothorax
  \item Tube thoracostomy (chest tube 24--32 Fr connected to water seal) for large or secondary pneumothorax
  \item Immediate needle decompression (14-gauge, second intercostal space, midclavicular line) for tension pneumothorax
  \item Pleurodesis (mechanical or chemical) or video-assisted thoracoscopic surgery (VATS) with blebectomy for recurrent pneumothorax
\end{itemize}

\section*{References}
\begin{thebibliography}{99}
\bibitem{pneumothorax_symptoms:ref1} Bense L, Wiman LG, Hedenstierna G. ``Spontaneous pneumothorax: a prospective study of 500 cases.'' \textit{Chest}, 1993;104(6):1769--1773.
\bibitem{pneumothorax_symptoms:ref2} MacDuff A, Arnold A, Harvey J, et al. ``Management of spontaneous pneumothorax: British Thoracic Society pleural disease guideline 2010.'' \textit{Thorax}, 2010;65(Suppl 2):ii18--ii31.
\bibitem{pneumothorax_symptoms:ref3} Baumann MH, Strange C, Heffner JE, et al. ``Management of spontaneous pneumothorax: an American College of Chest Physicians Delphi consensus statement.'' \textit{Chest}, 2001;119(2):590--602.
\bibitem{pneumothorax_symptoms:ref4} Noppen M, De Keukeleire T. ``Pneumothorax.'' \textit{Respiration}, 2008;76(2):121--127.
\bibitem{pneumothorax_symptoms:ref5} Sahn SA, Heffner JE. ``Spontaneous pneumothorax.'' \textit{New England Journal of Medicine}, 2000;342(12):868--874.
\bibitem{pneumothorax_symptoms:ref6} Gupta D, Hansell A, Nichols T, et al. ``Epidemiology of pneumothorax in England.'' \textit{Thorax}, 2000;55(8):666--671.

\end{thebibliography}
